\section{Estructura del ADN}

Los componentes básicos del ADN son los nucleótidos, los cuales se componen de lo siguiente:

\begin{itemize}
    \item Base
    \item Azúcar
    \item Grupo fosfato ($PO_{4}$)
\end{itemize}

\section{Niveles de estructura del ADN}

\begin{itemize}
    \item Primaria: Nivel definido por la secuancia de nucleótidos que hay en una hebra de ADN.
    
    \item Secundaria: Estructura de doble hélice formada por la unión de dos hebras que son:
	    \begin{itemize}
	        \item Complementarias: A cada base le corresponde otra base específica en la otra hebra.
	    	    \begin{itemize}    
	    	        \item A --> T
	    	        \item C --> G
	    	    \end{itemize}
	        \item Antiparalelas: Van en sentido opuesta la una de la otra.
	    \end{itemize}
    
    \item Terciaria: Las hélices de ADN se enrollan sobre proteínas llamadas histonas, las cuales compactan el espacio que ocupa el ADN

    \item Cuaternaria: Conjunto de nucleosomas llamado selenoide. Al darse la división celular, el selenoide sirve para compactar más el ADN en estructuras llamadas cromosomas.
\end{itemize}

\subsection{Histonas y nucleosomas}

Las proteínas anteriormente mencionadas, las histonas, son proteínas que enrrollan el ADN, y se componen de 5 partes:

\begin{itemize}
    \item H1: Eslabón o linker
    \item H2A, H2B, H3, H4: Core
\end{itemize}

Por otro lado, los nucleo somas son conjuntos de histonas que se agrupan entre sí.

\section{Proceso de replicación}

Es el proceso en que el ADN de una célula se replica para preparase para la división celular. Este proceso se compone de 3 partes:

\begin{enumerate}
    \item Iniciación: Adherencia de proteínas al sitio de replicación (Ori)
    \item Elongación: Aquí trabaja la ADN polimerasa, moviéndose de 5' a 3'
    \item Terminación: Actividad de endo/exonucleasa
\end{enumerate}

\subsection{Componentes de la replicación}

\begin{itemize}
    \item Topoisomerasa: Esta va  por fuera del bucle. Desenrolla la cadena de ADN para que, posteriormente, la helicasa pueda romperla.
    \item Helicasa: Esta va por dentro del bucle. Rompe los enlaces de hidrógeno que unen a la cadena de ADN.
    \item Fragmentos de Okasaki: Se componen de ARN.
    \item Lugar de unión: Lugar donde las proteínas se unen a la hebra rezagada.
    \item Nucleotidos: Se transforman de ADN a ARN por acción de la proteína ADN ligasa. 
        \begin{itemize}
            \item Procariotas: 1000-2000 nucleótidos.
            \item Eucariotas: Menor cantidad de nucleótidos.
        \end{itemize}
        
    \item Cadenas
        \begin{itemize}
            \item Cadena principal: Trabaja de 3' a 5'
            \item Cadena rezagada: Trabaja de 5' a 3'. Aquí se aplican los fragmentos de Okazaki
        \end{itemize}
    \item Proteínas de unión: Se unen a los puntos de iniciación
    \item Proteínas de unión monocatenarias: Bloquean que se vuelva a enrollar la cadena
    \item ADNpolimerasa (ADNpol): Trabaja de 5' a 3'. Necesita de un grupo hidroxilo libre ($OH^{-}$).
    \item Primers, cebadores o iniciadores: Se colocan por actuar de la ARN primasa. Proporcionan el grupo hidroxilo libre ($OH^{-}$) que la ADNpol necesita. Se componen de 5-10 nucleotidos.
\end{itemize}

\subsection{Primera parte: Iniciación}

Origen de replicación (Ori)

\subsection{Segunda parte: Elongación}



\subsection{Tercera parte: Terminación}



\section{Particularidades de la replicación eucariota}

Hay una serie de factores que influyen en que la replicación del ADN de células eucariotas sea algo más compleja. Estos son que el ADN de células eucariotas es más largo y se encuentra "enrredado" en proteínas llamadas histonas, las cuales al agregarse forman nucleosomas. Al darse la replicación del ADN, la hebra vieja se queda con los nucleosomas viejos y la hebra nueva se ``enrreda'' en histonas nuevas.